\section{The counterexample}\label{sec:construction}

We now construct the presentation complex and perfect normal subgroup
needed in Proposition~\ref{prop:plus}. The quotient must contain torsion
detected by a character, and the corresponding local system on the
presentation complex must have vanishing second homology.
Perfectness will follow from
annihilation relations with coprime coefficients in the subgroup's
abelianization.
Related uses of coprime annihilators for relation modules appear in
Bridson--Tweedale~\cite[Proposition~3.3]{BridsonTweedale07}.

Let $B$ be the presentation complex of
\begin{equation}\label{eq:presentation}
 P=\left\langle x_1,a_1,x_2,a_2,s\ \middle|\
 \begin{aligned}
 x_1a_1x_1^{-1}&=a_1^4,&
 x_2a_2x_2^{-1}&=a_2^3,\\
 x_1&=a_2^{13},&
 s x_2 s^{-1}&=a_1^5
 \end{aligned}\right\rangle.
\end{equation}
Thus $B$ has one vertex, five edges, and four two-dimensional cells.
Define
\[
 t_1=x_1^2,\qquad t_2=x_2^3,\qquad
 K=\langle\!\langle t_1,t_2\rangle\!\rangle_P,
 \qquad G=P/K,
\]
where the double brackets denote normal closure in $P$.
Sending $s$ to $1\in\Z$ and the other four generators to zero
respects all these relations and descends to an epimorphism
$G\to\Z$. In particular, $G$ is infinite.
We use the commutator convention $[g,h]=ghg^{-1}h^{-1}$.

\begin{lemma}\label{lem:perfect}
The normal subgroup $K$ is perfect.
\end{lemma}

\begin{proof}
The exponents are matched by $4^2-1=15=3\cdot5$ and
$3^3-1=26=2\cdot13$, with $\gcd(15,26)=1$.
Conjugation by $t_1=x_1^2$ sends $a_1$ to $a_1^{16}$, while
conjugation by $t_2=x_2^3$ sends $a_2$ to $a_2^{27}$.
Squaring $x_1=a_2^{13}$ and cubing $s x_2s^{-1}=a_1^5$ now give
\begin{equation}\label{eq:commutators}
 t_1=a_2^{26}=[t_2,a_2],
 \qquad
 s t_2 s^{-1}=a_1^{15}=[t_1,a_1].
\end{equation}
Write $K_{\ab}$ additively and let $u_j$ be the image of $t_j$.
Conjugation makes $K_{\ab}$ a left $\Z G$-module: every element
of $K$ acts trivially on its own abelianization. For $t\in K$,
the image of $[t,a]$ is $(1-a)[t]$. Consequently
\begin{equation}\label{eq:module-relations}
 u_1=(1-a_2)u_2,
 \qquad s u_2=(1-a_1)u_1.
\end{equation}

Since $t_1=a_2^{26}$, the element $a_2$ fixes $u_1$ and
$a_2^{26}$ acts trivially on all of $K_{\ab}$. Applying
$1+a_2+\cdots+a_2^{25}$ to the first equality gives
\[
 26u_1=(1-a_2^{26})u_2=0.
\]
Likewise, $a_1$ fixes $s u_2$, because $s t_2s^{-1}=a_1^{15}$;
and $a_1^{15}$ acts trivially on $K_{\ab}$. Applying
$1+a_1+\cdots+a_1^{14}$ to the second equality gives
\[
 15s u_2=(1-a_1^{15})u_1=0.
\]
The action of $s$ is invertible, so $15u_2=0$, and the first
equality of \eqref{eq:module-relations} then gives $15u_1=0$.
Since $15$ and $26$ are relatively prime, $u_1=0$. The second
equality gives $s u_2=0$, hence $u_2=0$.

Every element of $K$ is a product of $P$-conjugates of $t_1^{\pm1}$
and $t_2^{\pm1}$. Thus $u_1,u_2$ generate $K_{\ab}$ under
addition and the $P$-action. Their vanishing proves $K_{\ab}=0$.
\end{proof}

By Proposition~\ref{prop:plus}, there is a finite complex $X\supseteq B$
with fundamental group $G$, obtained by adding two two-dimensional
and two three-dimensional cells, such that
\begin{equation}\label{eq:d2-for-x}
 H_i(\widetilde X;\Z)=0\quad(i>2),\qquad
 H^3(X;M)=0\quad\text{for every left }\Z G\text{-module }M.
\end{equation}
It remains to find the character that obstructs a two-dimensional model.

Choose a primitive third root of unity $\zeta\in\C$. The assignments
\begin{equation}\label{eq:character}
 \rho(x_1)=-1,\quad \rho(a_1)=\zeta,\quad
 \rho(x_2)=\zeta^2,\quad \rho(a_2)=-1,\quad \rho(s)=1
\end{equation}
satisfy the four relations of \eqref{eq:presentation}, since
\[
 \zeta=\zeta^4,\qquad -1=(-1)^3,\qquad
 -1=(-1)^{13},\qquad \zeta^2=\zeta^5.
\]
They also send $t_1$ and $t_2$ to $1$. Hence they define a character
$\rho:G\to\C^\times$. Denote its local system on $X$ and its
pullback to $B$ by $L_\rho$.

\Needspace{6\baselineskip}
\begin{lemma}\label{lem:twisted-vanishing}
The local system $L_\rho$ satisfies $H_2(X;L_\rho)=0$.
\end{lemma}

\begin{proof}
Proposition~\ref{prop:plus} reduces the computation to $B$.
Its twisted cellular groups in degrees two and one are $\C^4$ and
$\C^5$, respectively. We compute the second boundary by the lifted
edge-path sum, equivalently by Fox differentiation~\cite{Fox53}.
In this convention the derivative of a word obeys
\[
 \frac{\partial(uv)}{\partial g}
   =\frac{\partial u}{\partial g}
      +u\frac{\partial v}{\partial g},\qquad
 \frac{\partial g}{\partial g}=1,\qquad
 \frac{\partial g^{-1}}{\partial g}=-g^{-1},
\]
and the derivatives with respect to other generators are zero.
Evaluate their group-ring coefficients at $\rho$.

Use the relators
\[
 \begin{split}
 r_1&=x_1a_1x_1^{-1}a_1^{-4},\qquad
 r_2=x_2a_2x_2^{-1}a_2^{-3},\\
 r_3&=x_1a_2^{-13},\qquad
 r_4=sx_2s^{-1}a_1^{-5}.
 \end{split}
\]
With the edges ordered as $x_1,a_1,x_2,a_2,s$, the four boundary
vectors, written as rows, are
\begin{equation}\label{eq:fox-matrix}
 D=\begin{pmatrix}
 1-\zeta & -2 & 0 & 0 & 0\\
 0&0&2&\zeta^2-1&0\\
 1&0&0&-1&0\\
 0&-(1+\zeta)&1&0&1-\zeta^2
 \end{pmatrix}.
\end{equation}
Here is a direct check of the entries. For a relator
$xax^{-1}a^{-n}$, the evaluated entries in the $x$ and $a$
coordinates are
\[
 1-\rho(a)^n,
 \qquad \rho(x)-\sum_{j=0}^{n-1}\rho(a)^j.
\]
For $r_1$, the latter sum is
$1+\zeta+\zeta^2+\zeta^3=1$; for $r_2$, it is $1-1+1=1$.
For $r_3$, the alternating sum of thirteen terms is $1$.
For $r_4$, the $a_1$ entry is
$-\sum_{j=0}^4\zeta^j=-(1+\zeta)$, the $x_2$ entry is
$\rho(s)=1$, and the $s$ entry is $1-\rho(x_2)=1-\zeta^2$.
These give \eqref{eq:fox-matrix}.

The minor using columns $x_1,a_1,x_2,s$ has determinant
\begin{equation}\label{eq:nonzero-minor}
 -4(1-\zeta^2)\ne0.
\end{equation}
Thus the four rows are independent, so the boundary
$C_2(B;L_\rho)\to C_1(B;L_\rho)$ is injective.
Since $B$ has no three-dimensional cells, $H_2(B;L_\rho)=0$,
as required.
\end{proof}

\begin{proof}[Proof of Theorem~\ref{thm:main}]
The complex $X$ constructed above is finite and satisfies
\eqref{eq:d2-for-x}. In $G$, the element $x_1$ has square one,
and it is nonidentity because $\rho(x_1)=-1$. Thus its order is
exactly two.

Suppose that $X$ were homotopy equivalent to a finite connected
complex $Y$ of dimension at most two. Transfer $\rho$ along the
induced isomorphism of fundamental groups, giving a local system
$L$ on $Y$. Homotopy invariance and Lemma~\ref{lem:twisted-vanishing}
give $H_2(Y;L)=0$. But the transferred character takes the value
$-1$ on an element of order two, contrary to
Corollary~\ref{cor:torsion}.

Concretely, a loop representing that element extends to a map
$\RP^2\to Y$: the attaching loop of the two-dimensional cell of
$\RP^2$ is its generator traversed twice. The pulled-back local
system has monodromy $-1$, whereas $b_1(\RP^2;\Q)=0$.
This is precisely the map excluded by Theorem~\ref{thm:tower}.
\end{proof}
